\chapter{How the floor moves with temperature}\label{ch:thermaln}

\section{The temperature dependence of a floor}\label{sec:floorT}

The left-hand mark of every device gauge is set by temperature alone: $\kB T_j\ln2$ at a chip's junction (Chapter~\ref{ch:cmos}) and, for a
qubit, the thermal floor of its record at the temperature the record sees (Chapter~\ref{ch:ascoreqc}). How that mark moves when the
temperature changes follows from the same expressions. For a record held in a gap $E=hf$, the held-record floor is
$p_{\rm eq}=1/(1+e^{M})$ with $M=hf/\kB T$, and its local logarithmic slope is
\begin{equation}
 \frac{d\ln p_{\rm eq}}{d\ln T}=M\,(1-p_{\rm eq})
 \label{eq:thermal_slope}
\end{equation}
\derived. The per-gate floor $p_{\rm th}=p_{\rm eq}t_g/T_1$ carries in addition the temperature dependence of $T_1$, so its slope is
Eq.~\eqref{eq:thermal_slope} minus $d\ln T_1/d\ln T$ \derived. Neither slope is a constant. For a transmon at 5~GHz the slope of $p_{\rm eq}$
is 6.85 at the qubit's own 35~mK and 16.0 at 15~mK \calc: the colder the record, the steeper its floor falls with further cooling.

The relaxation part has its own forms. For mechanisms with power-law densities of states the rate is a power of $T$, for example Raman
spin--lattice relaxation ($T^7$, $T^9$) or one-phonon processes ($T^1$ at high $T$)~\cite{Abragam1970}; for activated mechanisms
($e^{-E_a/\kB T}$) the logarithmic slope is $E_a/\kB T$ and varies with $T$ \derived. A single exponent in temperature therefore describes a
floor only over a stated window, and the window has to be named with the number.

\section{A thermal-emission benchmark}

Decoherence of C$_{70}$ matter waves by thermal emission~\cite{Hackermuller2004} shows the physics directly. Laser heating
before a Talbot--Lau interferometer raises the internal temperature from about 1000 to
3000~K. The resulting emission of visible photons carries which-path information and destroys
the fringes. The emission is thermally activated: with the radiative threshold $E_a\approx1.5$~eV the local slope $E_a/\kB T$ is 12.8 at
1360~K and 5.5 at 3140~K, and the secant across that window is 8.67 \calc.

\begin{checkbox}[title={Independent check: what the C$_{70}$ data can and cannot supply}]
\begin{itemize}
\item Figure~4 of the experiment plots visibility against \emph{heating laser power}. Its mean temperatures come from
      the model of heating and radiative cooling, and are not measured directly.
\item The emission rate of hot C$_{70}$ is thermally activated: it is negligible below about 2000~K, and the spectrum is cut off
      at long wavelengths by the lack of radiative transitions below $\sim1.5$~eV~\cite{Hackermuller2004}. Any exponent read over
      1360--3140~K is therefore a local slope of an activated process, and it would be different over a different window.
\item Decoherence by gas collisions is linear in gas \emph{pressure}~\cite{Hornberger2003}: one event per collision, linear in the
      density of the collision partners. At fixed pressure its temperature dependence goes as $n_{\mathrm{gas}}\bar v\propto T^{-1/2}$,
      and at fixed density as $T^{+1/2}$ \derived. It does not supply an exponent of 1 in temperature.
\end{itemize}
The C$_{70}$ experiment shows that thermal emission converts a quantum object into a
classical one, which is exactly the physics IAM is about. It shows equally that the temperature dependence of a decoherence
channel is a property of that channel over a window, to be measured on the device in question.
\end{checkbox}

\begin{remark}[Which temperature, in kelvin]
A floor depends on an absolute temperature. Halving a Celsius reading does not halve the temperature and gives no meaningful factor \calc. The
temperature must also be the one the noise source sees. For a transmon it is the effective temperature of the qubit
mode or of the quasiparticles, which commonly lies above the mixing-chamber reading. For ions and neutral atoms it is
the motional temperature after laser cooling, set by the cooling limit and the heating rate, not by any cryostat
stage; the Doppler limit of Yb$^+$ on its 369~nm line is $\hbar\Gamma/2\kB=0.47$~mK for $\Gamma/2\pi=19.6$~MHz
\calc. For NV centres the spin relaxation rate near room temperature follows a two-phonon Raman $T^5$ term plus an
Orbach term, not a single power~\cite{Jarmola2012}.
For a neutral-atom Rydberg gate two temperatures matter: the atoms' motion, near 10~$\mu$K, and the radiation field around the apparatus, which at
300~K places 62--620 thermal photons in each mode of a 100--10~GHz Rydberg transition and drives transitions out of the Rydberg state
\calc~\cite{Beterov2009}. Chapter~\ref{ch:qplatforms} gives the temperature that matters for each platform.
\end{remark}

\section{Why the slope matters in practice}
Cooling moves the floor, and Eq.~\eqref{eq:thermal_slope} says by how much. Halving a transmon's own temperature from 35 to 17.5~mK at 5~GHz
raises $M$ from 6.86 to 13.7 and lowers $p_{\rm eq}$ from $1.05\times10^{-3}$ to $1.1\times10^{-6}$ \calc. Because present two-qubit errors sit
far above the per-gate thermal floor (Chapter~\ref{ch:ascoreqc}), what cooling buys a gate is decided by the channels that are themselves
thermal, such as the qubit's residual excited population at the start of a gate and thermally activated quasiparticles (Chapter~\ref{ch:xqp}),
and each of those is measured against the refrigerator on the same device \interp. Figure~\ref{fig:qp_floors} (Chapter~\ref{ch:qplatforms})
draws the thermal occupation of every platform's gap against the temperature its record sees. The measurement that settles how far cooling
pays for a given processor is the one Ref.~\cite{Jin2015} made: the excited-state occupation and $T_1$ against refrigerator temperature, with
the two-qubit error measured at each point \openprob.
